\section{Background and Related Work}

\subsection{Tcl's Core Data Model}

Tcl's \texttt{Tcl\_Obj} structure maintains both a string representation and an internal
representation, with reference counting for automatic memory management and copy-on-write
(COW) for shared references. Lists and dictionaries are highly optimized built-in types
that exploit this infrastructure. VOO leverages these existing mechanisms rather than
reimplementing them.

\subsection{Itcl and TclOO}

Itcl introduced OO to Tcl as a separate system layered atop the language, resulting in
fundamental incompatibilities: its \texttt{variable} command conflicted with Tcl's
namespace \texttt{variable}; its access control was bypassable through namespace
manipulation; class redefinition was blocked, violating Tcl's dynamic nature; and
multiple incompatible destruction mechanisms existed. These challenges originated from
designing OO constructs independently of Tcl's philosophical principles.

TclOO (TIP~\#257)~\cite{fellows2006} was a deliberate departure from Itcl, with a new
\texttt{oo::} namespace, a minimalist core, improved namespace alignment, and unified
destruction via \texttt{destroy}. However, persistent limitations remain: objects are
handles requiring explicit destruction, omitted \texttt{destroy} calls cause memory
leaks, each object maintains a procedure handle with computational and memory overhead,
and converting lists/dicts to TclOO objects requires API changes.

\subsection{Lessons from Two Decades}

The evolution reveals that philosophical consistency is critical, clean architectural
breaks can be justified, extensibility enables ecosystem growth, and migration paths
are essential. VOO applies these lessons by maximizing philosophical alignment,
providing incremental migration paths, and prioritizing simplicity over comprehensive
features.
